\section{Campos finitos y ejemplos de cálculo mediante discriminantes}\label{sec:campos-finitos}

\subsection{Frobenius y el ejemplo de sesenta y cuatro elementos}\label{tema:27-2-2}
Sea $F$ un campo finito de característica $p$. Contiene al campo primo $\mathbb F_p$. Si $d=[F:\mathbb F_p]$, una base identifica su conjunto subyacente con $\mathbb F_p^d$, luego $|F|=p^d$. Para cada $a\ne0$, Lagrange aplicado a $F^\times$ da $a^{p^d-1}=1$, y por tanto todos los elementos de $F$, incluido cero, satisfacen
\[
a^{p^d}-a=0.
\]
El polinomio $t^{p^d}-t$ tiene derivada $-1$, así que es separable. Sus $p^d$ raíces son exactamente los elementos de $F$, que es su campo de descomposición sobre $\mathbb F_p$. En particular, $F/\mathbb F_p$ es de Galois. La unicidad del campo de descomposición muestra que dos campos con $p^d$ elementos son isomorfos; se utiliza la notación $\mathbb F_{p^d}$.

También existe un campo con ese número de elementos para cada $d\geq1$. En un campo de descomposición de $t^{p^d}-t$, el conjunto de sus raíces es estable bajo suma, producto, opuestos e inversos de elementos no nulos, porque la potencia $p^d$ es un homomorfismo de campos. Las raíces forman un subcampo que contiene a $\mathbb F_p$ y genera el campo de descomposición; por ello son todo el campo. Como la derivada es $-1$, hay exactamente $p^d$ raíces distintas.

El automorfismo de Frobenius
\[
\varphi:\mathbb F_{p^d}\longrightarrow\mathbb F_{p^d},
\qquad a\longmapsto a^p
\]
fija a $\mathbb F_p$ y satisface $\varphi^d=\id$. Su orden es exactamente $d$: si $0<r<d$ y $\varphi^r=\id$, los $p^d$ elementos del campo serían raíces de $t^{p^r}-t$, cuyo grado $p^r$ es menor. Por el conteo de automorfismos de una extensión de Galois,
\[
\gal(\mathbb F_{p^d}/\mathbb F_p)=\ang{\varphi}\cong C_d.
\]

\begin{ejem}\label{ejem:C27-f64}
Para $\mathbb F_{64}/\mathbb F_2$, el grupo es $C_6$. Sus únicos subgrupos tienen órdenes $1,2,3,6$, y corresponden respectivamente a campos de tamaños $64,8,4,2$. Así los únicos campos intermedios son $\mathbb F_{64}$, $\mathbb F_8$, $\mathbb F_4$ y $\mathbb F_2$, con el par de retículos de la \cref{fig:reticulo-f64}:
\figuraApuntes{Figuras/ReticuloF64}{fig:reticulo-f64}{Correspondencia de Galois para $\mathbb F_{64}/\mathbb F_2$.}
En particular,
\[
\Aut_{\mathbb F_8}(\mathbb F_{64})=\ang{\varphi^3}=\{\id,\varphi^3\},
\qquad [\mathbb F_{64}:\mathbb F_8]=2.
\]
Análogamente, $\Aut_{\mathbb F_4}(\mathbb F_{64})=\ang{\varphi^2}$ tiene orden tres. El grado de cada campo sobre $\mathbb F_2$ es el índice del subgrupo correspondiente en $C_6$.
\end{ejem}

\subsection{La acción sobre las raíces}\label{tema:27-3-1}
Para un polinomio separable $f\in k[t]$ con campo de descomposición $F/k$, escribimos
\[
\gal_k(f)=\gal(F/k).
\]
La elección de otro campo de descomposición produce un grupo isomorfo, mediante conjugación por un isomorfismo sobre $k$.

\begin{lema}\label{lema:C27-transitiva}
Si $f\in k[t]$ es separable e irreducible de grado $n$, su grupo de Galois actúa transitivamente sobre las raíces. Al numerarlas se identifica con un subgrupo transitivo de $S_n$.
\end{lema}
\begin{proof}
Todo automorfismo sobre $k$ permuta las raíces de $f$. La acción es fiel: si fija todas las raíces, fija el campo que generan, que es $F$.

Sean $\alpha$ y $\beta$ dos raíces. Por irreducibilidad existe un isomorfismo $k(\alpha)\to k(\beta)$ que fija a $k$ y envía $\alpha$ a $\beta$. El campo $F$ es de descomposición de $f$ sobre ambos campos. La unicidad con transporte del campo base de la \cref{prop:C24-descomposicion} prolonga ese isomorfismo a un automorfismo de $F$. Por ello cualquier raíz puede enviarse a cualquier otra.
\end{proof}

En grado tres, los únicos subgrupos transitivos, salvo conjugación, son $A_3$ y $S_3$: su orden es múltiplo de tres por órbita-estabilizador, y divide a seis por Lagrange.

En grado cuatro, las posibilidades son
\[
S_4,\qquad A_4,\qquad D_8,\qquad C_2\times C_2,\qquad C_4.
\]
Para justificar la lista, el orden de un subgrupo transitivo de $S_4$ es múltiplo de cuatro y divide a $24$, luego es $4,8,12$ o $24$. Un grupo de orden cuatro es $C_4$ o $C_2\times C_2$, y su acción transitiva sobre cuatro puntos es regular: los estabilizadores son triviales. Los subgrupos de orden ocho son los de Sylow de $S_4$, conjugados al grupo de simetrías de un cuadrado, $D_8$. Un subgrupo de orden doce tiene índice dos y es el núcleo de un homomorfismo sobreyectivo $S_4\to C_2$. Todas las transposiciones son conjugadas y generan $S_4$, de modo que todas se envían al elemento no trivial; el homomorfismo es el signo y su núcleo es $A_4$. El caso de orden $24$ es $S_4$.

\subsection{El discriminante y el grupo alternante}\label{tema:27-3-2}
Sea $f\in k[t]$ mónico y separable, con raíces $\alpha_1,\ldots,\alpha_n$ en su campo de descomposición. Definimos
\[
\Delta=\prod_{1\leq i<j\leq n}(\alpha_i-\alpha_j),
\qquad D=\Delta^2.
\]
El elemento $D$ es el \emph{discriminante} de $f$. Como las raíces son distintas, $\Delta\ne0$. Una transposición de dos raíces cambia el signo de $\Delta$; por composición, cualquier permutación $\pi$ satisface $\pi(\Delta)=\operatorname{sgn}(\pi)\Delta$. Por tanto toda permutación fija a $D$, y la correspondencia de Galois da $D\in k$. Para estudiar un polinomio no mónico mediante este criterio, podemos dividirlo por su coeficiente principal: el asociado mónico tiene las mismas raíces y el mismo grupo de Galois.

\begin{lema}\label{lema:C27-discriminante}
Supongamos que $\operatorname{char}k\ne2$. Para un polinomio mónico separable $f$ de grado $n$, se tiene $\gal_k(f)\subseteq A_n$ si y sólo si su discriminante es un cuadrado en $k$.
\end{lema}
\begin{proof}
Como $\Delta\ne0$ y $-\Delta\ne\Delta$, el grupo fija a $\Delta$ exactamente cuando sus permutaciones son pares. Por la condición de campo fijo, esto equivale a $\Delta\in k$, que hace de $D=\Delta^2$ un cuadrado del campo base.

Si $D=d^2$ para algún $d\in k$, entonces $\Delta^2=d^2$, de donde $(\Delta-d)(\Delta+d)=0$. Como estamos en un campo, $\Delta=d$ o $\Delta=-d$. Así $\Delta\in k$ y todas las permutaciones del grupo son pares.
\end{proof}

Para una cúbica mónica
\[
f(t)=t^3+at^2+bt+c,
\]
la expansión del producto de diferencias, usando las relaciones entre raíces y coeficientes, da
\[
D=a^2b^2-4a^3c-4b^3+18abc-27c^2.
\]
Si además la característica es distinta de tres, conviene eliminar el término cuadrático:
\[
f\left(t-\frac a3\right)=t^3+pt+q,
\qquad p=b-\frac{a^2}{3},\qquad q=c-\frac{ab}{3}+\frac{2a^3}{27}.
\]
Trasladar todas las raíces por la misma cantidad no cambia sus diferencias, por lo que no cambia $D$. Para $g(t)=t^3+pt+q$, la fórmula se reduce a
\[
D=-4p^3-27q^2.
\]
Podemos comprobarla directamente. Si $r,s,u$ son las raíces, entonces
\[
r+s+u=0,\quad rs+ru+su=p,\quad rsu=-q,
\]
de donde $r^2+s^2+u^2=-2p$ y $r^2s^2+r^2u^2+s^2u^2=p^2$. Además,
\[
\begin{aligned}
-D&=g'(r)g'(s)g'(u)\\
&=(3r^2+p)(3s^2+p)(3u^2+p)\\
&=27q^2+9p\,p^2+3p^2(-2p)+p^3
=27q^2+4p^3.
\end{aligned}
\]
La primera igualdad sigue de $g'(r)=(r-s)(r-u)$ y las dos identidades análogas. Sustituir $p$ y $q$ recupera la fórmula general. Esta última es una identidad con coeficientes enteros, por lo que también vale en característica tres, aunque allí no se puede efectuar el cambio $t\mapsto t-a/3$.

Si la cúbica es irreducible y separable y la característica es distinta de dos, su grupo es $A_3$ cuando $D$ es un cuadrado, y $S_3$ cuando no lo es: ésas son las dos posibilidades transitivas del grado tres.

\subsection{Ejemplos cúbicos y normalidad}\label{tema:27-3-3}\label{tema:27-4-0b}
\begin{ejem}[El campo de descomposición de $t^3-2$]\label{ejem:C27-s3}
El polinomio $t^3-2$ es irreducible sobre $\mathbb Q$ por Eisenstein con $2$ y es separable por característica cero. Su discriminante es
\[
D=-27(-2)^2=-108,
\]
que no es un cuadrado racional. Por tanto $\gal_{\mathbb Q}(t^3-2)\cong S_3$.

Pongamos $a=\sqrt[3]{2}\in\mathbb R$ y $\omega=e^{2\pi i/3}$. El campo de descomposición es
\[
L=\mathbb Q(a,a\omega,a\omega^2)=\mathbb Q(a,\omega).
\]
La primera igualdad con la segunda se debe a que $\omega=(a\omega)/a$ y a que las tres raíces pertenecen a $\mathbb Q(a,\omega)$. Como $\mathbb Q(a)\subseteq\mathbb R$ y $\omega$ no es real, su mínimo sobre ese campo es $t^2+t+1$. Así $[L:\mathbb Q]=2\cdot3=6$, de acuerdo con el orden de $S_3$.

Podemos elegir los automorfismos
\[
\begin{array}{c|cc}
&a&\omega\\\hline
\sigma&a\omega&\omega\\
\tau&a&\omega^2
\end{array}
\]
Como $[L:\mathbb Q(\omega)]=3$, el mínimo de $a$ sobre $\mathbb Q(\omega)$ es $t^3-2$. La existencia de $\sigma$ se obtiene extendiendo el isomorfismo que envía $a$ a $a\omega$ sobre esa base, mediante la unicidad de campos de descomposición. El automorfismo $\tau$ es la conjugación compleja. Se cumple
\[
\sigma^3=\tau^2=\id,\qquad \tau\sigma\tau=\sigma^{-1},
\qquad \gal(L/\mathbb Q)=\ang{\sigma,\tau}\cong S_3.
\]
La relación de conjugación se comprueba en los generadores $a$ y $\omega$; los seis elementos $\sigma^i$ y $\tau\sigma^i$, $0\leq i<3$, son distintos por sus imágenes de esos generadores.

Los dos retículos se muestran en la \cref{fig:reticulo-cubica}.
\figuraApuntes{Figuras/ReticuloCubica}{fig:reticulo-cubica}{Correspondencia de Galois para $\mathbb Q(\sqrt[3]{2},\omega)/\mathbb Q$.}
El campo $\mathbb Q(\omega)$ corresponde a $\ang{\sigma}=A_3$; su grado sobre $\mathbb Q$ es dos y $[L:\mathbb Q(\omega)]=3$. Los tres campos cúbicos corresponden a los tres subgrupos de orden dos, que tienen índice tres. En concreto,
\[
\begin{aligned}
L^{\ang{\tau}}&=\mathbb Q(a),\\
L^{\ang{\tau\sigma}}&=\mathbb Q(a\omega),\\
L^{\ang{\tau\sigma^2}}&=\mathbb Q(a\omega^2).
\end{aligned}
\]
Con la composición de funciones, el subgrupo indicado fija al generador correspondiente. El campo fijo y el campo de la derecha tienen grado tres sobre $\mathbb Q$, por lo que coinciden.

La conjugación por $\sigma$ permuta los tres subgrupos de orden dos y sus campos fijos. Por ejemplo, $\tau$ fija a $a$, pero $\sigma$ envía $\mathbb Q(a)$ a $\mathbb Q(a\omega)\ne\mathbb Q(a)$. Por tanto $\ang{\tau}$ no es normal en $S_3$, y $\mathbb Q(a)/\mathbb Q$ no es normal. En cambio, las extensiones de $L$ sobre todos sus campos intermedios sí son de Galois.
\end{ejem}

\begin{ejem}[Un grupo cíclico de orden tres]
El polinomio $f(t)=t^3-3t+1$ no tiene raíces racionales: las únicas posibilidades son $1$ y $-1$, y ninguna lo anula. Por ser cúbico, es irreducible sobre $\mathbb Q$. Es separable y su discriminante es
\[
D=-4(-3)^3-27=108-27=81=9^2.
\]
Así $\gal_{\mathbb Q}(f)=A_3\cong C_3$. Su campo de descomposición tiene grado tres sobre $\mathbb Q$; en particular, queda generado por cualquiera de sus raíces, porque cada una tiene grado tres y pertenece a él.
\end{ejem}

Para el campo de descomposición de $t^3-2$ sobre $\mathbb Q$, $A_3$ es normal en $S_3$ y
\[
\gal(\mathbb Q(\omega)/\mathbb Q)\cong S_3/A_3\cong C_2.
\]
Los tres subgrupos de orden dos no son normales; sus campos cúbicos correspondientes son separables sobre $\mathbb Q$, pero no son extensiones de Galois.
